\section{Introduction}
\label{introduction}
% Providing background information on the topic at the beginning would greatly benefit readers, as it helps establish context and highlight the significance of the research.
In the span of 20 years, the Armenian economy has experienced four major crises and one positive economic upturn, all of which have been triggered by external shocks. Even though all of these are major events, the specific risks and propagation channels behind them are fundamentally different. In this paper, we quantify the systemic risks relevant to the Armenian financial system, differentiate between various types of risks, and compare the buildup and materialization of these risks before and during the major crises. Additionally, we assess the impact of macroprudential policy on these risks and propose a macroprudential policy rule that considers both the risks and resilience of the Armenian financial system.

We distinguish between two types of systemic risks: financial vulnerability and financial stress. Financial vulnerability refers to structural issues in the system that accumulate over time and have the potential to amplify economic shocks during crises. To measure financial vulnerability, we use the Financial Cycle Index (FCyI) \citep{hollo2012, plasil2014}. In contrast, financial stress describes a specific state of the system, characterized by heightened market volatility, tight credit conditions and liquidity constraints. To measure financial stress, we rely on the financial conditions index (FCI) \citep{koop_2013}. Finally, to estimate the effects of macroprudential policy, the current literature \citep{alam2019, galan2020} suggests constructing a macroprudential policy index (MPI) based on the history and the macroprudential stance of a country.

However, as in many countries, the data related to macroprudential policy in Armenian is quite limited. To address this, we employ a three step approach. First, we estimate a panel quantile regression model, following a methodology similar to that of \citet{galan2020} and \citet{aikman2019} as a robustness check. We derive results for the effects of macroprudential policy on the distribution of output growth. This type of inference has been popularized in recent years and is a major part of the growth-at-risk (GaR) framework \citep{esrb2024}. We focus on output growth because it is closely related to the performance of the financial system and serves as a useful measure for quantifying vulnerabilities and potential losses.

Next, we estimate a similar, but local quantile regression model tailored specifically to Armenia, including the limited macroprudential data. Finally, we compare the results from the panel model with the macroprudential policy estimates from the local model. This comparison helps us identify the generalized relationships captured by the panel model, which we can use as a source of informative priors to compensate for the lack of policy data in the local model.

According to the results, systemic risks play a major role in shaping the GDP growth distribution. Financial vulnerabilities negatively affect the lower tail of the future growth distribution, while their impact on the median is either negative or insignificant. An increase in the FCyI indicates higher vulnerabilities; this relationship is best characterized by a leveraged bubble \citep{jorda2015}, that may emerge from positive feedback loops between credit growth and asset prices, and upon collapse create a severe recession with slower recovery. Meanwhile, an increase in FCI reflects higher financial stress. The deterioration of the financial conditions has a negative effect on current GDP growth, while its effect on future growth is either positive or insignificant.

Macroprudential policy exhibits nonlinear effects on the short-term growth distribution. The effects are negative on the median of the distribution. For instance, when the countercyclical capital buffer (CCyB) is tightened, stricter capital requirements are imposed on banks, leading to a reduction in the credit supply. This, in turn, leads to a decrease in consumption and investment, thereby negatively affecting short-term output growth. However, our findings also show that macroprudential policy has positive effects on the tail of the distribution. Continuing with the example for CCyB, higher capital enhances the resilience of the banks and reduces potential losses in the event of tail risks. \citet{jimnez2017} also find that an increase in capital buffers extends credit to firms and reduces the risk of insolvency. Our results are consistent with the existing literature \citep{galan2020} and support the theoretical framework where macroprudential policy is facing a trade-off between short-term costs (a decline in output) and long-term benefits (reduction of risks captured by the shrinking tail of the distribution).

Using this framework, it is possible to perform conditional forecasts for the median and the tail of the GDP growth distribution. We also refer to the difference between the median and tail as “distance-to-tail” or “Gap.” The conditional distance-to-tail serves as a measure of the overall systemic risk in the economy stemming from financial exposures. By tightening the MPI, we reduce the median (lower reward), but the distance-to-tail is also reduced (lower risk). The benefits of this trade-off for policymakers will depend on their definition of welfare.

In our paper we consider two district welfare functions to arrive at the optimal policy rule: the distance-to-tail approach \citep{esrb2024} and the \citet{suarez2022} welfare function. The distance-to-tail rule focuses on systemic risks, comparing the estimated distance-to-tail values with systemic thresholds derived from the quantiles of the historical series. The thresholds serve as an implicit measure of the resilience of the financial system. When the estimated distance-to-tail exceeds these thresholds, it indicates that systemic risks are underestimated, warranting a tightening of macroprudential policy. Meanwhile, the utility function is based on the welfare approach developed by \citet{suarez2022}. For the policy rule, \citet{suarez2022} proposes a simple welfare function derived from microfoundations, incorporating both the median and distance-to-tail of the GDP growth distribution. This allows us to capture the trade-offs that macroprudential policy imposes. If the long-term benefits outweigh the short-term costs, then the policy should be tightened. We use our empirical framework alongside this utility function to determine the “optimal policy.” 

The difference between the welfare functions is subtle, yet crucial. The distance-to-tail rule considers the trade-off between risk and resilience, without accounting for the potential costs of the policy. In contrast, the welfare function proposed by \citet{suarez2022} emphasizes the trade-off between the costs and benefits of the policy, while overlooking resilience. Neither approach is inherently superior; rather, they are complementary, providing insights with both conservative and opportunistic perspectives on risk management.

% The second-to-last paragraph of the introduction provides a brief overview of the main contributions. While all three contributions are relevant from a policy perspective, the first two (1. Assessment of systemic risk in Armenia; 2. Use of panel and localized models) are not particularly strong from an academic standpoint. We suggest emphasizing the specific advantages and strengths of these contributions to better showcase their value.
This paper makes several contributions to both academic research and policy development. First, it provides a comprehensive assessment of systemic risks in the Armenian financial system, offering insights into the specific vulnerabilities faced by small, open economies that are particularly sensitive to external shocks. Second, it employs both panel and localized models to measure the impact of macroprudential policy, examining how different approaches can affect the results. Finally, the paper presents a flexible framework and a set of indicators that are instrumental for macroprudential policy decision-making and can be applied to other emerging economies. We have performed this analysis with a strong intention to use it in the policymaking procedure.

The rest of the paper has the following structure. Section 2 provides an overview of systemic risks and the Growth-at-Risk (GaR) framework. Section 3 evaluates the role of macroprudential policy and its effects on both risks and resilience. Section 4 introduces the methodological approach and describes the data used in the study. We discuss and interpret the quantile regression results in Section 5. Section 6 will discuss the two approaches for the macroprudential policy stance and their relevance for making informed policy decisions. Finally, Section 7 offers concluding remarks and discusses the broader policy implications of the framework.